\documentclass[11pt,letterpaper]{article}
\usepackage[margin=0.8in]{geometry}
\usepackage[T1]{fontenc}
\usepackage[utf8]{inputenc}
\usepackage{lmodern,microtype,xcolor,enumitem,hyperref,fancyhdr,needspace}
\definecolor{ink}{HTML}{183CB5}
\hypersetup{colorlinks=true,urlcolor=ink,linkcolor=ink,pdftitle={CS4308 Graph Machine Learning Reading List},pdfauthor={CS4308 Course Materials}}
\pagestyle{fancy}\fancyhf{}\fancyhead[L]{\small CS4308 / Graph Machine Learning}\fancyhead[R]{\small Reading companion}\fancyfoot[R]{\thepage}
\setlength{\headheight}{14pt}\setlength{\parindent}{0pt}\setlength{\parskip}{5pt}
\setlist[itemize]{leftmargin=*,nosep}\urlstyle{same}
\begin{document}
{\Large\bfseries CS4308 Graph Machine Learning}\par
{\huge\bfseries Fifty papers to read and build}\par
\vspace{5pt}
A curated path from graph structure and random walks to neural message passing, transformers, generation, and reliable evaluation. The selection combines the instructor's graph-learning research with foundational methods, evaluation frameworks, and selected 2024--2026 research from the wider field. Bibliographic sources were checked in November 2026. The year shown is the cited publication year; an arXiv preprint may be earlier. Workshop papers are identified explicitly.
\section*{How to use this list}
The class collectively discusses all 50 papers. Each group, normally a pair, must understand, implement, evaluate, and present at least five distinct papers. Register a balanced selection: one random-walk or spectral method, one message-passing or attention method, one advanced method, one evaluation or reliability paper, and one project-aligned paper. For theory and benchmark papers, implement a central diagnostic or evaluation protocol.

Submit a critical brief, reproducible code, a central result, one ablation or stress test, and a presentation for each study. Small-scale reproductions are acceptable when explicitly scoped. Both partners must contribute and defend the work. Suggested checkpoints are weeks 4, 6, 9, 11, and 13. This reading companion uses the same P01--P50 identifiers as the course outline and website.

GAN means generative adversarial network; GAT means graph attention network. The original GAN paper supplies the general prerequisite, followed by GraphGAN and MolGAN for graph-based adversarial learning.
\section*{Weekly reading map}

\textbf{Week 2: Spectral structure and graph descriptors.} P01, P06\par
\textbf{Week 3: Walk embeddings and statistical representations.} P03, P04, P05, P07, P08\par
\textbf{Week 4: From graph signals to GCNs.} P12, P13, P16\par
\textbf{Week 5: Message passing and attention.} P14, P15, P17, P20\par
\textbf{Week 6: Graph prediction and feature construction.} P10, P19, P21, P22\par
\textbf{Week 7: Expressivity, failure modes and evaluation.} P25, P26, P48\par
\textbf{Week 9: Scaling, self-supervision and transfer.} P09, P11, P18, P27, P28, P30, P31\par
\textbf{Week 10: Graph learning in practice.} P02, P32, P33, P34, P35, P36, P37, P38\par
\textbf{Week 11: Graph transformers and tokenization.} P23, P39, P40, P41, P44\par
\textbf{Week 12: Adversarial and generative graph learning.} P29, P42, P43, P45\par
\textbf{Week 13: Trustworthy graph research.} P46, P47, P49, P50\par
\textbf{Week 14: Research symposium and foundation models.} P24\par
\clearpage\section*{Selection and coverage}
The 50-paper set includes 22 publications coauthored by Mudassir Shabbir, 19 field landmarks and nine recent research additions. Seventeen readings were published in 2024-2026: eight in 2024, six in 2025 and three in 2026.\par
Author coverage comes first: include the publicly verifiable graph-learning papers, graph descriptors, graph-learning applications, and directly relevant graph-distinguishability foundations coauthored by Mudassir Shabbir. Coverage was checked on 2 November 2026; this does not assert coverage of unpublished or unlisted manuscripts.\par
The 28 wider-field readings balance 19 foundational methods and evaluation frameworks with nine recent papers: three additions each from 2024, 2025 and 2026. Selection prioritizes methodological significance, influence and course relevance. Recent work is labeled as such; especially for 2026, long-term impact is not yet established. PageRank, DeepWalk, node2vec, GCN, GAN, GAT and graph transformers remain covered.\par
A preprint and its published version count as one reading. Distinct conference and extended journal papers are retained where they make separate contributions; their relationship is noted. Preprints and works in progress are labeled. General geometry and control papers without a direct graph-learning contribution are outside this reading set.\par
All nine recent additions have verified official conference proceedings records. Years follow the published venue, not the first arXiv posting: for example, GraphAny is an ICLR 2025 paper. Combined with the instructor's papers, the list contains eight publications from 2024, six from 2025 and three from 2026.\par
All 50 receive class coverage through lectures, concise paper spotlights and discussions. The assigned week identifies the main discussion; groups may choose papers across weeks for their five-paper portfolio.\par
Coverage sources: \href{https://orcid.org/0000-0002-6961-0961}{author ORCID}; \href{https://anwar-said.github.io/anwarsaid/CV.pdf}{coauthor publication list}; \href{https://amessbee.github.io/s27/reading_selection.json}{coverage and version audit}. Each entry links its primary publication or manuscript.\par
\clearpage
\Needspace{16\baselineskip}\section*{Foundations and representations}
\begin{minipage}{\linewidth}
{\large\bfseries P01 The PageRank Citation Ranking: Bringing Order to the Web}\par
Lawrence Page et al. Stanford InfoLab Technical Report 1999-66, 1999.\par
\textcolor{ink}{\small FIELD LANDMARK}\par
\textbf{Why read it.} Connects graph centrality, Markov chains, and ranking before learned representations.\par
\textbf{Build or investigate.} Implement power iteration with teleportation; test dangling nodes and sensitivity to the damping factor.\par
\textbf{Level.} Foundation\quad\href{http://ilpubs.stanford.edu:8090/422/}{Primary source $\nearrow$}\par
\end{minipage}\par\vspace{10pt}
\begin{minipage}{\linewidth}
{\large\bfseries P03 DeepWalk: Online Learning of Social Representations}\par
Bryan Perozzi et al. KDD, 2014.\par
\textcolor{ink}{\small FIELD LANDMARK}\par
\textbf{Why read it.} Turns short random walks into a practical unsupervised representation-learning pipeline.\par
\textbf{Build or investigate.} Implement walk generation and skip-gram training; evaluate node classification with a fixed split.\par
\textbf{Level.} Foundation\quad\href{https://arxiv.org/abs/1403.6652}{Primary source $\nearrow$}\par
\end{minipage}\par\vspace{10pt}
\begin{minipage}{\linewidth}
{\large\bfseries P04 node2vec: Scalable Feature Learning for Networks}\par
Aditya Grover and Jure Leskovec. KDD, 2016.\par
\textcolor{ink}{\small FIELD LANDMARK}\par
\textbf{Why read it.} Shows how biased walks trade off local neighborhoods and structural roles.\par
\textbf{Build or investigate.} Sweep return and exploration parameters against DeepWalk using identical downstream evaluation.\par
\textbf{Level.} Foundation\quad\href{https://arxiv.org/abs/1607.00653}{Primary source $\nearrow$}\par
\end{minipage}\par\vspace{10pt}
\begin{minipage}{\linewidth}
{\large\bfseries P05 Estimating Descriptors for Large Graphs}\par
Zohair Raza Hassan, Mudassir Shabbir, Imdadullah Khan, Waseem Abbas. PAKDD 2020, LNCS 12084, 779-791, 2020.\par
\textcolor{ink}{\small MUDASSIR SHABBIR / RESEARCH}\par
{\small Conference paper; the extended edge-stream journal paper is a separate reading.}\par
\textbf{Why read it.} Streaming graph descriptors for classification.\par
\textbf{Build or investigate.} Implement reservoir-sampled graphlet counts; plot descriptor error and classification accuracy against the memory budget.\par
\textbf{Level.} Intermediate\quad\href{https://doi.org/10.1007/978-3-030-47426-3_60}{Primary source $\nearrow$}\par
\end{minipage}\par\vspace{10pt}
\begin{minipage}{\linewidth}
{\large\bfseries P06 Interpretable Multi-Scale Graph Descriptors via Structural Compression}\par
Ammar Ahmed, Zohair Raza Hassan, Mudassir Shabbir. Information Sciences 533, 169-180, 2020.\par
\textcolor{ink}{\small MUDASSIR SHABBIR / RESEARCH}\par
\textbf{Why read it.} Interpretable multiscale graph representations for classification via structural compression.\par
\textbf{Build or investigate.} Implement structural compression at several scales; compare descriptor discrimination and runtime with degree and spectral summaries.\par
\textbf{Level.} Intermediate\quad\href{https://doi.org/10.1016/j.ins.2020.05.032}{Primary source $\nearrow$}\par
\end{minipage}\par\vspace{10pt}
\begin{minipage}{\linewidth}
{\large\bfseries P07 DGSD: Distributed Graph Representation via Graph Statistical Properties}\par
Anwar Said, Saeed-Ul Hassan, Suppawong Tuarob, Raheel Nawaz, Mudassir Shabbir. Future Generation Computer Systems 119, 166-175, 2021.\par
\textcolor{ink}{\small MUDASSIR SHABBIR / RESEARCH}\par
\textbf{Why read it.} Distributed permutation-invariant graph embedding.\par
\textbf{Build or investigate.} Construct statistical descriptors under different graph partitions; test permutation invariance and compare with a degree-histogram baseline.\par
\textbf{Level.} Intermediate\quad\href{https://doi.org/10.1016/j.future.2021.02.005}{Primary source $\nearrow$}\par
\end{minipage}\par\vspace{10pt}
\begin{minipage}{\linewidth}
{\large\bfseries P08 NetKI: A Kirchhoff Index Based Statistical Graph Embedding in Nearly Linear Time}\par
Anwar Said, Saeed-Ul Hassan, Waseem Abbas, Mudassir Shabbir. Neurocomputing 433, 108-118, 2021.\par
\textcolor{ink}{\small MUDASSIR SHABBIR / RESEARCH}\par
\textbf{Why read it.} Scalable graph representation and graph classification using effective resistance.\par
\textbf{Build or investigate.} Compute exact effective resistances on small graphs, compare an approximate NetKI representation, and measure runtime versus classification quality.\par
\textbf{Level.} Intermediate\quad\href{https://doi.org/10.1016/j.neucom.2020.12.075}{Primary source $\nearrow$}\par
\end{minipage}\par\vspace{10pt}
\begin{minipage}{\linewidth}
{\large\bfseries P09 Computing Graph Descriptors on Edge Streams}\par
Zohair Raza Hassan, Sarwan Ali, Imdadullah Khan, Mudassir Shabbir, Waseem Abbas. ACM Transactions on Knowledge Discovery from Data 17(8), article 112, 2023.\par
\textcolor{ink}{\small MUDASSIR SHABBIR / RESEARCH}\par
{\small Journal extension of the 2020 descriptor paper; reproduce its additional contribution.}\par
\textbf{Why read it.} Streaming structural/spectral graph descriptors and graph classification.\par
\textbf{Build or investigate.} Reproduce one streaming descriptor and compare against exact computation across memory budgets and stream orderings; isolate the journal extension.\par
\textbf{Level.} Intermediate\quad\href{https://doi.org/10.1145/3591468}{Primary source $\nearrow$}\par
\end{minipage}\par\vspace{10pt}
\begin{minipage}{\linewidth}
{\large\bfseries P10 On Augmenting Topological Graph Representations for Attributed Graphs}\par
Anwar Said, Mudassir Shabbir, Saeed-Ul Hassan, Zohair Raza Hassan, Ammar Ahmed, Xenofon Koutsoukos. Applied Soft Computing 136, 110104, 2023.\par
\textcolor{ink}{\small MUDASSIR SHABBIR / RESEARCH}\par
\textbf{Why read it.} Integrating graph attributes with topological representations for graph classification.\par
\textbf{Build or investigate.} Compare topology-only, attributes-only and augmented representations under fixed splits to test when node attributes add useful information.\par
\textbf{Level.} Intermediate\quad\href{https://doi.org/10.1016/j.asoc.2023.110104}{Primary source $\nearrow$}\par
\end{minipage}\par\vspace{10pt}
\begin{minipage}{\linewidth}
{\large\bfseries P11 MSDGSD: A Scalable Graph Descriptor for Processing Large Graphs}\par
Muhammad Ali, Anwar Said, Iqra Safder, Saeed Ul Hassan, Naif Radi Aljohani, Mudassir Shabbir. IEEE Transactions on Computational Social Systems 11(3), 3594-3605, 2024.\par
\textcolor{ink}{\small MUDASSIR SHABBIR / RESEARCH}\par
\textbf{Why read it.} Scalable graph descriptors for downstream machine learning.\par
\textbf{Build or investigate.} Implement a tractable multiscale descriptor; compare DGSD and MSDGSD at matched compute and measure accuracy and runtime as graph size grows.\par
\textbf{Level.} Intermediate\quad\href{https://doi.org/10.1109/TCSS.2023.3338691}{Primary source $\nearrow$}\par
\end{minipage}\par\vspace{10pt}
\Needspace{16\baselineskip}\section*{Graph neural networks}
\begin{minipage}{\linewidth}
{\large\bfseries P12 The Graph Neural Network Model}\par
Franco Scarselli, Marco Gori, Ah Chung Tsoi, Markus Hagenbuchner, Gabriele Monfardini. IEEE Transactions on Neural Networks 20(1), 61-80, 2009.\par
\textcolor{ink}{\small FIELD LANDMARK}\par
\textbf{Why read it.} An early general formulation of neural computation on graphs; provides historical grounding for today's message-passing architectures.\par
\textbf{Build or investigate.} Implement a small recurrent graph-state update and examine convergence; compare its assumptions with a fixed-depth GCN.\par
\textbf{Level.} Advanced\quad\href{https://doi.org/10.1109/TNN.2008.2005605}{Primary source $\nearrow$}\par
\end{minipage}\par\vspace{10pt}
\begin{minipage}{\linewidth}
{\large\bfseries P13 Convolutional Neural Networks on Graphs with Fast Localized Spectral Filtering}\par
Michaël Defferrard et al. NeurIPS (then NIPS), 2016.\par
\textcolor{ink}{\small FIELD LANDMARK}\par
\textbf{Why read it.} Replaces expensive eigendecomposition with localized polynomial filters.\par
\textbf{Build or investigate.} Implement Chebyshev recurrence; compare polynomial order, receptive field, runtime, and accuracy.\par
\textbf{Level.} Intermediate\quad\href{https://arxiv.org/abs/1606.09375}{Primary source $\nearrow$}\par
\end{minipage}\par\vspace{10pt}
\begin{minipage}{\linewidth}
{\large\bfseries P14 Inductive Representation Learning on Large Graphs}\par
William L. Hamilton et al. NeurIPS (then NIPS), 2017.\par
\textcolor{ink}{\small FIELD LANDMARK}\par
\textbf{Why read it.} Introduces neighborhood sampling and aggregation for generalization to unseen nodes.\par
\textbf{Build or investigate.} Implement mean aggregation with sampling and evaluate a genuine inductive node split.\par
\textbf{Level.} Intermediate\quad\href{https://arxiv.org/abs/1706.02216}{Primary source $\nearrow$}\par
\end{minipage}\par\vspace{10pt}
\begin{minipage}{\linewidth}
{\large\bfseries P15 Neural Message Passing for Quantum Chemistry}\par
Justin Gilmer et al. ICML, 2017.\par
\textcolor{ink}{\small FIELD LANDMARK}\par
\textbf{Why read it.} Unifies many graph architectures through message, update, and readout functions.\par
\textbf{Build or investigate.} Build a small edge-aware message-passing model for one molecular regression target.\par
\textbf{Level.} Intermediate\quad\href{https://proceedings.mlr.press/v70/gilmer17a.html}{Primary source $\nearrow$}\par
\end{minipage}\par\vspace{10pt}
\begin{minipage}{\linewidth}
{\large\bfseries P16 Semi-Supervised Classification with Graph Convolutional Networks}\par
Thomas N. Kipf and Max Welling. ICLR, 2017.\par
\textcolor{ink}{\small FIELD LANDMARK}\par
\textbf{Why read it.} Establishes a compact, reproducible baseline for learning from node features and edges.\par
\textbf{Build or investigate.} Implement a two-layer GCN; ablate self-loops, normalization, and feature-only baselines.\par
\textbf{Level.} Foundation\quad\href{https://arxiv.org/abs/1609.02907}{Primary source $\nearrow$}\par
\end{minipage}\par\vspace{10pt}
\begin{minipage}{\linewidth}
{\large\bfseries P17 Graph Attention Networks}\par
Petar Veli\v{c}kovi\'c et al. ICLR, 2018.\par
\textcolor{ink}{\small FIELD LANDMARK}\par
\textbf{Why read it.} Introduces learned neighborhood weighting and multi-head attention.\par
\textbf{Build or investigate.} Implement masked graph attention; compare heads and attention dropout with GCN.\par
\textbf{Level.} Intermediate\quad\href{https://arxiv.org/abs/1710.10903}{Primary source $\nearrow$}\par
\end{minipage}\par\vspace{10pt}
\begin{minipage}{\linewidth}
{\large\bfseries P19 How Powerful are Graph Neural Networks?}\par
Keyulu Xu et al. ICLR, 2019.\par
\textcolor{ink}{\small FIELD LANDMARK}\par
\textbf{Why read it.} Relates aggregation to graph isomorphism and motivates the Graph Isomorphism Network.\par
\textbf{Build or investigate.} Implement GIN and test sum versus mean aggregation on controlled multiset and graph examples.\par
\textbf{Level.} Advanced\quad\href{https://arxiv.org/abs/1810.00826}{Primary source $\nearrow$}\par
\end{minipage}\par\vspace{10pt}
\begin{minipage}{\linewidth}
{\large\bfseries P20 Enhanced Graph Neural Networks with Ego-Centric Spectral Subgraph Embeddings Augmentation}\par
Anwar Said, Mudassir Shabbir, Tyler Derr, Waseem Abbas, Xenofon Koutsoukos. 22nd IEEE International Conference on Machine Learning and Applications (ICMLA), 2023.\par
\textcolor{ink}{\small MUDASSIR SHABBIR / RESEARCH}\par
\textbf{Why read it.} Spectral ego-subgraph embeddings augment GNN node features.\par
\textbf{Build or investigate.} Augment a GCN with ego-subgraph spectral features; vary neighborhood radius and include feature-only and random-feature controls.\par
\textbf{Level.} Intermediate\quad\href{https://doi.org/10.1109/ICMLA58977.2023.00341}{Primary source $\nearrow$}\par
\end{minipage}\par\vspace{10pt}
\begin{minipage}{\linewidth}
{\large\bfseries P21 Improving Graph Machine Learning Performance Through Feature Augmentation Based on Network Control Theory}\par
Anwar Said, Obaid Ullah Ahmad, Waseem Abbas, Mudassir Shabbir, Xenofon Koutsoukos. 32nd Mediterranean Conference on Control and Automation (MED), 322-327, 2024.\par
\textcolor{ink}{\small MUDASSIR SHABBIR / RESEARCH}\par
{\small Conference study; compare with the expanded 2025 rank-encoding paper.}\par
\textbf{Why read it.} Network control based node-feature augmentation for GNNs.\par
\textbf{Build or investigate.} Add network-control features to a fixed GNN; compare no features, random features and spectral features while holding the training budget constant.\par
\textbf{Level.} Intermediate\quad\href{https://doi.org/10.1109/MED61351.2024.10566241}{Primary source $\nearrow$}\par
\end{minipage}\par\vspace{10pt}
\begin{minipage}{\linewidth}
{\large\bfseries P22 Network Controllability Perspectives on Graph Representation}\par
Anwar Said, Obaid Ullah Ahmad, Waseem Abbas, Mudassir Shabbir, Xenofon Koutsoukos. IEEE Transactions on Knowledge and Data Engineering 36(8), 4116-4128, 2024.\par
\textcolor{ink}{\small MUDASSIR SHABBIR / RESEARCH}\par
\textbf{Why read it.} Control-theoretic graph representation for machine learning.\par
\textbf{Build or investigate.} Construct a control-theoretic graph representation and compare it with spectral and degree features under the same classifier and split.\par
\textbf{Level.} Advanced\quad\href{https://doi.org/10.1109/TKDE.2023.3331318}{Primary source $\nearrow$}\par
\end{minipage}\par\vspace{10pt}
\begin{minipage}{\linewidth}
{\large\bfseries P23 Feature Construction Using Network Control Theory and Rank Encoding for Graph Machine Learning}\par
Anwar Said, Yifan Wei, Obaid Ullah Ahmad, Mudassir Shabbir, Waseem Abbas, Xenofon Koutsoukos. IEEE Open Journal of Control Systems, 2025.\par
\textcolor{ink}{\small MUDASSIR SHABBIR / RESEARCH}\par
{\small Expanded journal study following the 2024 feature-augmentation paper.}\par
\textbf{Why read it.} Controllability and rank encoding generate expressive node features for graph classification.\par
\textbf{Build or investigate.} Reproduce rank-based node features on small benchmarks; compare the 2024 feature construction, rank encoding and standard positional encodings.\par
\textbf{Level.} Advanced\quad\href{https://doi.org/10.1109/OJCSYS.2025.3599371}{Primary source $\nearrow$}\par
\end{minipage}\par\vspace{10pt}
\Needspace{16\baselineskip}\section*{Expressivity and limits}
\begin{minipage}{\linewidth}
{\large\bfseries P25 On the Bottleneck of Graph Neural Networks and its Practical Implications}\par
Uri Alon and Eran Yahav. ICLR, 2021.\par
\textcolor{ink}{\small FIELD LANDMARK}\par
\textbf{Why read it.} Introduces over-squashing as a failure of information transmission across graph bottlenecks.\par
\textbf{Build or investigate.} Reproduce a small long-range task and compare added shortcuts with increased hidden dimension.\par
\textbf{Level.} Advanced\quad\href{https://arxiv.org/abs/2006.05205}{Primary source $\nearrow$}\par
\end{minipage}\par\vspace{10pt}
\begin{minipage}{\linewidth}
{\large\bfseries P26 Control-based Conditions for Graph Distinguishability}\par
Muhammad Raza, Obaid Ullah Ahmad, Mudassir Shabbir, Xenofon Koutsoukos, Waseem Abbas. SIAM Conference on Control and its Applications (CT), 31-38, 2025.\par
\textcolor{ink}{\small MUDASSIR SHABBIR / RESEARCH}\par
\textbf{Why read it.} Uses control-theoretic conditions to study graph distinguishability, connecting spectral structure with the limits of graph representations.\par
\textbf{Build or investigate.} Enumerate small graph pairs and test walk-matrix conditions alongside spectral and 1-WL signatures; report what each criterion can and cannot distinguish.\par
\textbf{Level.} Advanced\quad\href{https://doi.org/10.1137/1.9781611978742.5}{Primary source $\nearrow$}\par
\end{minipage}\par\vspace{10pt}
\Needspace{16\baselineskip}\section*{Scaling and self-supervision}
\begin{minipage}{\linewidth}
{\large\bfseries P28 Graph Contrastive Learning with Augmentations}\par
Yuning You et al. NeurIPS, 2020.\par
\textcolor{ink}{\small FIELD LANDMARK}\par
\textbf{Why read it.} Makes graph augmentations a central design choice in contrastive pretraining.\par
\textbf{Build or investigate.} Implement two graph augmentations and assess whether they preserve labels on the chosen dataset.\par
\textbf{Level.} Intermediate\quad\href{https://arxiv.org/abs/2010.13902}{Primary source $\nearrow$}\par
\end{minipage}\par\vspace{10pt}
\begin{minipage}{\linewidth}
{\large\bfseries P30 Control-Based Graph Embeddings with Data Augmentation for Contrastive Learning}\par
Obaid Ullah Ahmad, Anwar Said, Mudassir Shabbir, Xenofon Koutsoukos, Waseem Abbas. American Control Conference (ACC), 27-32, 2024.\par
\textcolor{ink}{\small MUDASSIR SHABBIR / RESEARCH}\par
\textbf{Why read it.} Unsupervised contrastive graph representation with control-informed augmentation.\par
\textbf{Build or investigate.} Compare control-informed graph augmentations with random edge deletion under the same contrastive objective; ablate augmentation strength.\par
\textbf{Level.} Advanced\quad\href{https://doi.org/10.23919/ACC60939.2024.10644876}{Primary source $\nearrow$}\par
\end{minipage}\par\vspace{10pt}
\begin{minipage}{\linewidth}
{\large\bfseries P31 Learning Backbones: Sparsifying Graphs Through Zero Forcing for Effective Graph-Based Learning}\par
Obaid Ullah Ahmad, Anwar Said, Mudassir Shabbir, Xenofon Koutsoukos, Waseem Abbas. Complex Networks \& Their Applications XIII (COMPLEX NETWORKS 2024), 85-100, 2025.\par
\textcolor{ink}{\small MUDASSIR SHABBIR / RESEARCH}\par
{\small Conference held in 2024; proceedings published in 2025. The arXiv version is the same reading.}\par
\textbf{Why read it.} Zero-forcing graph sparsification for downstream GNN classification.\par
\textbf{Build or investigate.} Extract zero-forcing backbones; compare with random and degree-based edge retention at matched sparsity, reporting accuracy, runtime and connectivity.\par
\textbf{Level.} Intermediate\quad\href{https://doi.org/10.1007/978-3-031-82427-2_8}{Primary source $\nearrow$}\par
\end{minipage}\par\vspace{10pt}
\Needspace{16\baselineskip}\section*{Applications and graph optimization}
\begin{minipage}{\linewidth}
{\large\bfseries P32 Leveraging Deep Learning and SNA Approaches for Smart City Policing in the Developing World}\par
Saeed-Ul Hassan, Mudassir Shabbir, Sehrish Iqbal, Anwar Said, Faisal Kamiran, Raheel Nawaz, Umar Saif. International Journal of Information Management 56, 102045, 2021.\par
\textcolor{ink}{\small MUDASSIR SHABBIR / RESEARCH}\par
\textbf{Why read it.} Direct GCN application to binary node classification in communication networks; suitable critical discussion of privacy, bias and evaluation.\par
\textbf{Build or investigate.} Use synthetic communication graphs to study label prediction and structural bias; evaluate leakage and false-positive disparities without profiling real individuals.\par
\textbf{Level.} Intermediate\quad\href{https://doi.org/10.1016/j.ijinfomgt.2019.102045}{Primary source $\nearrow$}\par
\end{minipage}\par\vspace{10pt}
\begin{minipage}{\linewidth}
{\large\bfseries P33 Attack-Resilient Multi-Agent Flocking Control Using Graph Neural Networks}\par
Chandreyee Bhowmick, Mudassir Shabbir, Xenofon Koutsoukos. 30th Mediterranean Conference on Control and Automation (MED), 300-305, 2022.\par
\textcolor{ink}{\small MUDASSIR SHABBIR / RESEARCH}\par
\textbf{Why read it.} GCN-based distributed flocking resilient to adversarial communication.\par
\textbf{Build or investigate.} Simulate a small flock under compromised neighbor messages; compare the GNN controller with a nominal controller using collision and cohesion metrics.\par
\textbf{Level.} Intermediate\quad\href{https://cs.vuse.vanderbilt.edu/koutsoxd/www/Publications/Attack-Resilient_Multi-Agent_Flocking_Control_Using_Graph_Neural_Networks.pdf}{Primary source $\nearrow$}\par
\end{minipage}\par\vspace{10pt}
\begin{minipage}{\linewidth}
{\large\bfseries P34 Circuit Design Completion Using Graph Neural Networks}\par
Anwar Said, Mudassir Shabbir, Brian Broll, Waseem Abbas, Peter V\"olgyesi, Xenofon Koutsoukos. Neural Computing and Applications 35, 12145-12157, 2023.\par
\textcolor{ink}{\small MUDASSIR SHABBIR / RESEARCH}\par
\textbf{Why read it.} GNNs for missing circuit components and connections.\par
\textbf{Build or investigate.} Mask a component or connection in small circuit graphs; compare GNN completion with frequency and topology baselines using circuit-level held-out splits.\par
\textbf{Level.} Intermediate\quad\href{https://doi.org/10.1007/s00521-023-08346-x}{Primary source $\nearrow$}\par
\end{minipage}\par\vspace{10pt}
\begin{minipage}{\linewidth}
{\large\bfseries P35 Learning-Based Heuristic for Combinatorial Optimization of the Minimum Dominating Set Problem using Graph Convolutional Networks}\par
Abihith Kothapalli, Mudassir Shabbir, Xenofon Koutsoukos. arXiv:2306.03434, 2023.\par
\textcolor{ink}{\small MUDASSIR SHABBIR / RESEARCH}\par
{\small Preprint.}\par
\textbf{Why read it.} GCNs for NP-hard graph combinatorial optimization.\par
\textbf{Build or investigate.} Train on small synthetic graphs; compare feasible dominating sets against greedy and exact solutions, then test size and graph-family generalization.\par
\textbf{Level.} Intermediate\quad\href{https://arxiv.org/abs/2306.03434}{Primary source $\nearrow$}\par
\end{minipage}\par\vspace{10pt}
\begin{minipage}{\linewidth}
{\large\bfseries P36 Sequential Graph Neural Networks for Source Code Vulnerability Identification}\par
Ammar Ahmed, Anwar Said, Mudassir Shabbir, Xenofon Koutsoukos. HotSoS 2023 work in progress; arXiv:2306.05375, 2023.\par
\textcolor{ink}{\small MUDASSIR SHABBIR / RESEARCH}\par
{\small Work-in-progress/preprint version; no journal publication is claimed.}\par
\textbf{Why read it.} Sequential/GNN code representation, vulnerability classification and CVEFGE dataset.\par
\textbf{Build or investigate.} Build a small code-graph pipeline and compare GNN versus sequence-only features; split by project and remove near-duplicate functions before evaluation.\par
\textbf{Level.} Intermediate\quad\href{https://arxiv.org/abs/2306.05375}{Primary source $\nearrow$}\par
\end{minipage}\par\vspace{10pt}
\begin{minipage}{\linewidth}
{\large\bfseries P37 A Graph Machine Learning Framework to Compute Zero Forcing Sets in Graphs}\par
Obaid Ullah Ahmad, Mudassir Shabbir, Waseem Abbas, Xenofon Koutsoukos. IEEE Transactions on Network Science and Engineering 11(2), 2110-2123, 2024.\par
\textcolor{ink}{\small MUDASSIR SHABBIR / RESEARCH}\par
\textbf{Why read it.} Graph ML solution to zero-forcing set computation.\par
\textbf{Build or investigate.} Train a graph learner to propose zero-forcing sets; verify feasibility by simulation and compare set size and runtime with greedy and exact small-graph solvers.\par
\textbf{Level.} Advanced\quad\href{https://doi.org/10.1109/TNSE.2023.3337750}{Primary source $\nearrow$}\par
\end{minipage}\par\vspace{10pt}
\Needspace{16\baselineskip}\section*{Relational and temporal learning}
\begin{minipage}{\linewidth}
{\large\bfseries P02 Towards Foundation Models for Knowledge Graph Reasoning}\par
Mikhail Galkin, Xinyu Yuan, Hesham Mostafa, Jian Tang, Zhaocheng Zhu. ICLR 2024, 2024.\par
\textcolor{ink}{\small RECENT RESEARCH / 2024}\par
\textbf{Why read it.} A significant step toward transferable knowledge-graph reasoning across unseen entity and relation vocabularies; connects graph representations with foundation-model generalization.\par
\textbf{Build or investigate.} Use an available pretrained model on two small held-out knowledge graphs; compare zero-shot inference and limited fine-tuning with a task-specific baseline. Audit graph overlap.\par
\textbf{Level.} Advanced\quad\href{https://proceedings.iclr.cc/paper_files/paper/2024/hash/85dd09d356ca561169b2c03e43cf305e-Abstract-Conference.html}{Primary source $\nearrow$}\par
\end{minipage}\par\vspace{10pt}
\Needspace{16\baselineskip}\section*{Graph foundation models}
\begin{minipage}{\linewidth}
{\large\bfseries P27 One For All: Towards Training One Graph Model For All Classification Tasks}\par
Hao Liu, Jiarui Feng, Lecheng Kong, Ningyue Liang, Dacheng Tao, Yixin Chen, Muhan Zhang. ICLR 2024, 2024.\par
\textcolor{ink}{\small RECENT RESEARCH / 2024}\par
\textbf{Why read it.} Studies a unified representation for node, link and graph tasks using text attributes and graph prompting; a useful reference for cross-domain graph learning.\par
\textbf{Build or investigate.} Reproduce a small two-domain transfer experiment using cached text embeddings; ablate prompt substructures and compare against separately trained GNNs.\par
\textbf{Level.} Advanced\quad\href{https://proceedings.iclr.cc/paper_files/paper/2024/hash/57faf5642eb06e0602b95f6aa989b38a-Abstract-Conference.html}{Primary source $\nearrow$}\par
\end{minipage}\par\vspace{10pt}
\begin{minipage}{\linewidth}
{\large\bfseries P18 Fully-inductive Node Classification on Arbitrary Graphs}\par
Jianan Zhao, Zhaocheng Zhu, Mikhail Galkin, Hesham Mostafa, Michael Bronstein, Jian Tang. ICLR 2025, 2025.\par
\textcolor{ink}{\small RECENT RESEARCH / 2025}\par
\textbf{Why read it.} Extends inductive learning to graphs with different feature and label spaces; provides a sharp experimental test of what cross-graph generalization means.\par
\textbf{Build or investigate.} Evaluate a released GraphAny model on two unseen small graphs with differing feature dimensions; compare against linear propagation and a supervised GCN without target-test tuning.\par
\textbf{Level.} Advanced\quad\href{https://proceedings.iclr.cc/paper_files/paper/2025/hash/3bdeb28a531f7af94b56bcdf8ee88f17-Abstract-Conference.html}{Primary source $\nearrow$}\par
\end{minipage}\par\vspace{10pt}
\begin{minipage}{\linewidth}
{\large\bfseries P24 Towards Graph Foundation Models: Learning Generalities Across Graphs via Task-Trees}\par
Zehong Wang, Zheyuan Zhang, Tianyi Ma, Nitesh V. Chawla, Chuxu Zhang, Yanfang Ye. ICML 2025, PMLR 267, 65518-65555, 2025.\par
\textcolor{ink}{\small RECENT RESEARCH / 2025}\par
\textbf{Why read it.} Uses task-trees to align node-, edge- and graph-level learning; connects transferability arguments with cross-task and cross-domain experiments.\par
\textbf{Build or investigate.} Build task-trees for one small dataset and two prediction tasks; test a frozen or lightly fine-tuned encoder against a task-specific baseline and audit pretraining overlap.\par
\textbf{Level.} Advanced\quad\href{https://proceedings.mlr.press/v267/wang25eq.html}{Primary source $\nearrow$}\par
\end{minipage}\par\vspace{10pt}
\Needspace{16\baselineskip}\section*{Graph transformers}
\begin{minipage}{\linewidth}
{\large\bfseries P40 Do Transformers Really Perform Bad for Graph Representation?}\par
Chengxuan Ying et al. NeurIPS, 2021.\par
\textcolor{ink}{\small FIELD LANDMARK}\par
\textbf{Why read it.} Shows how structural encodings make global attention useful for graph prediction.\par
\textbf{Build or investigate.} Implement a compact Graphormer and ablate degree and shortest-path encodings.\par
\textbf{Level.} Advanced\quad\href{https://arxiv.org/abs/2106.05234}{Primary source $\nearrow$}\par
\end{minipage}\par\vspace{10pt}
\begin{minipage}{\linewidth}
{\large\bfseries P41 Recipe for a General, Powerful, Scalable Graph Transformer}\par
Ladislav Ramp\'ašek et al. NeurIPS, 2022.\par
\textcolor{ink}{\small FIELD LANDMARK}\par
\textbf{Why read it.} Combines local message passing, global attention, and position or structure encodings.\par
\textbf{Build or investigate.} Implement one GPS block and perform a three-component ablation under a fixed parameter budget.\par
\textbf{Level.} Advanced\quad\href{https://arxiv.org/abs/2205.12454}{Primary source $\nearrow$}\par
\end{minipage}\par\vspace{10pt}
\begin{minipage}{\linewidth}
{\large\bfseries P39 Relational Graph Transformer}\par
Vijay Prakash Dwivedi, Sri Jaladi, Yangyi Shen, Federico Lopez, Charilaos Kanatsoulis, Rishi Puri, Matthias Fey, Jure Leskovec. ICLR 2026, 2026.\par
\textcolor{ink}{\small RECENT RESEARCH / 2026}\par
\textbf{Why read it.} A recent transformer design for heterogeneous temporal graphs derived from relational tables; examines how tokenization and local/global attention adapt to database structure.\par
\textbf{Build or investigate.} Compare RelGT and a GNN on one small RelBench task at matched compute; ablate temporal or structural token components and report memory and predictive performance.\par
\textbf{Level.} Advanced\quad\href{https://proceedings.iclr.cc/paper_files/paper/2026/hash/fa1ef43326cfb8a808c0649708187f2a-Abstract-Conference.html}{Primary source $\nearrow$}\par
\end{minipage}\par\vspace{10pt}
\begin{minipage}{\linewidth}
{\large\bfseries P44 Graph Tokenization for Bridging Graphs and Transformers}\par
Zeyuan Guo, Enmao Diao, Cheng Yang, Chuan Shi. ICLR 2026, 2026.\par
\textcolor{ink}{\small RECENT RESEARCH / 2026}\par
\textbf{Why read it.} Explores reversible graph serialization and substructure-guided tokenization for standard sequence transformers, offering a fresh alternative to specialized graph architectures.\par
\textbf{Build or investigate.} Serialize small graphs, train or reuse a compact tokenizer, and compare a small transformer with a GNN; measure sensitivity to node permutations, sequence length and vocabulary size.\par
\textbf{Level.} Advanced\quad\href{https://proceedings.iclr.cc/paper_files/paper/2026/hash/2c0781f2ed2d8054d4f27a19bd8c380b-Abstract-Conference.html}{Primary source $\nearrow$}\par
\end{minipage}\par\vspace{10pt}
\Needspace{16\baselineskip}\section*{Generative graph learning}
\begin{minipage}{\linewidth}
{\large\bfseries P42 Generative Adversarial Nets}\par
Ian J. Goodfellow et al. NeurIPS (then NIPS), 2014.\par
\textcolor{ink}{\small FIELD LANDMARK}\par
\textbf{Why read it.} Provides the adversarial-learning prerequisite; this paper itself is not a graph model.\par
\textbf{Build or investigate.} Implement a small GAN on a toy distribution and diagnose instability before attempting graph generation.\par
\textbf{Level.} Foundation\quad\href{https://papers.nips.cc/paper/2014/hash/f033ed80deb0234979a61f95710dbe25-Abstract.html}{Primary source $\nearrow$}\par
\end{minipage}\par\vspace{10pt}
\begin{minipage}{\linewidth}
{\large\bfseries P43 GraphGAN: Graph Representation Learning with Generative Adversarial Nets}\par
Hongwei Wang et al. AAAI, 2018.\par
\textcolor{ink}{\small FIELD LANDMARK}\par
\textbf{Why read it.} Applies adversarial training to connectivity distributions for graph embeddings.\par
\textbf{Build or investigate.} Implement a small generator-discriminator link model and compare with a random-walk embedding baseline.\par
\textbf{Level.} Advanced\quad\href{https://arxiv.org/abs/1711.08267}{Primary source $\nearrow$}\par
\end{minipage}\par\vspace{10pt}
\begin{minipage}{\linewidth}
{\large\bfseries P45 MolGAN: An implicit generative model for small molecular graphs}\par
Nicola De Cao and Thomas Kipf. ICML Workshop on Theoretical Foundations and Applications of Deep Generative Models, 2018.\par
\textcolor{ink}{\small FIELD LANDMARK}\par
\textbf{Why read it.} Combines adversarial graph generation with property objectives and exposes mode collapse.\par
\textbf{Build or investigate.} Implement a reduced small-molecule generator; report validity, uniqueness, novelty, and diversity.\par
\textbf{Level.} Advanced\quad\href{https://arxiv.org/abs/1805.11973}{Primary source $\nearrow$}\par
\end{minipage}\par\vspace{10pt}
\begin{minipage}{\linewidth}
{\large\bfseries P29 A Large-scale Training Paradigm for Graph Generative Models}\par
Yu Wang, Ryan Rossi, Namyong Park, Huiyuan Chen, Nesreen Ahmed, Puja Trivedi, Franck Dernoncourt, Danai Koutra, Tyler Derr. ICLR 2025, 2025.\par
\textcolor{ink}{\small RECENT RESEARCH / 2025}\par
\textbf{Why read it.} Studies graph generation across multiple domains and text conditioning, moving the discussion from single-dataset fitting to reusable generative models.\par
\textbf{Build or investigate.} Evaluate a released checkpoint on a small held-out graph family; compare zero-shot and fine-tuned generation using validity, diversity, structural statistics and inference cost.\par
\textbf{Level.} Advanced\quad\href{https://proceedings.iclr.cc/paper_files/paper/2025/hash/3a2ef31a1e45908901adc0ca853a8faf-Abstract-Conference.html}{Primary source $\nearrow$}\par
\end{minipage}\par\vspace{10pt}
\Needspace{16\baselineskip}\section*{Evaluation and trust}
\begin{minipage}{\linewidth}
{\large\bfseries P47 GNNExplainer: Generating Explanations for Graph Neural Networks}\par
Rex Ying et al. NeurIPS, 2019.\par
\textcolor{ink}{\small FIELD LANDMARK}\par
\textbf{Why read it.} Explores which subgraphs and features support a GNN prediction.\par
\textbf{Build or investigate.} Implement explanation masks on a motif dataset; measure fidelity, sparsity, and stability.\par
\textbf{Level.} Intermediate\quad\href{https://arxiv.org/abs/1903.03894}{Primary source $\nearrow$}\par
\end{minipage}\par\vspace{10pt}
\begin{minipage}{\linewidth}
{\large\bfseries P48 Open Graph Benchmark: Datasets for Machine Learning on Graphs}\par
Weihua Hu et al. NeurIPS, 2020.\par
\textcolor{ink}{\small FIELD LANDMARK}\par
\textbf{Why read it.} Supplies task-specific datasets, realistic splits, and standardized evaluators.\par
\textbf{Build or investigate.} Implement an OGB baseline and document why its official split and metric match the task.\par
\textbf{Level.} Foundation\quad\href{https://arxiv.org/abs/2005.00687}{Primary source $\nearrow$}\par
\end{minipage}\par\vspace{10pt}
\begin{minipage}{\linewidth}
{\large\bfseries P49 A Survey of Graph Unlearning}\par
Anwar Said, Ngoc N. Tran, Yuying Zhao, Tyler Derr, Mudassir Shabbir, Waseem Abbas, Xenofon Koutsoukos. arXiv:2310.02164; revised October 2025, 2023.\par
\textcolor{ink}{\small MUDASSIR SHABBIR / RESEARCH}\par
{\small First posted in 2023; this reading uses the October 2025 revision.}\par
\textbf{Why read it.} Survey of graph machine unlearning, privacy, evaluation and applications.\par
\textbf{Build or investigate.} Select one surveyed unlearning method; compare deletion utility, retained-task performance and compute against full retraining using a clearly stated threat model.\par
\textbf{Level.} Intermediate\quad\href{https://arxiv.org/abs/2310.02164}{Primary source $\nearrow$}\par
\end{minipage}\par\vspace{10pt}
\begin{minipage}{\linewidth}
{\large\bfseries P50 NeuroGraph: Benchmarks for Graph Machine Learning in Brain Connectomics}\par
Anwar Said, Roza G. Bayrak, Tyler Derr, Mudassir Shabbir, Daniel Moyer, Catie Chang, Xenofon Koutsoukos. NeurIPS 2023 Datasets and Benchmarks, 2023.\par
\textcolor{ink}{\small MUDASSIR SHABBIR / RESEARCH}\par
\textbf{Why read it.} Graph ML benchmark for brain connectomics.\par
\textbf{Build or investigate.} Reproduce a small benchmark with authorized data, subject-disjoint splits and a non-graph baseline; test graph-construction sensitivity without making clinical claims.\par
\textbf{Level.} Intermediate\quad\href{https://arxiv.org/abs/2306.06202}{Primary source $\nearrow$}\par
\end{minipage}\par\vspace{10pt}
\begin{minipage}{\linewidth}
{\large\bfseries P38 RelBench: A Benchmark for Deep Learning on Relational Databases}\par
Joshua Robinson, Rishabh Ranjan, Weihua Hu, Kexin Huang, Jiaqi Han, Alejandro Dobles, Matthias Fey, Jan E. Lenssen, Yiwen Yuan, Zecheng Zhang, Xinwei He, Jure Leskovec. NeurIPS 2024 Datasets and Benchmarks, 2024.\par
\textcolor{ink}{\small RECENT RESEARCH / 2024}\par
\textbf{Why read it.} Establishes a benchmark for graph-based prediction over linked database tables, making temporal leakage and comparisons with tabular feature engineering central research questions.\par
\textbf{Build or investigate.} Run one small relational prediction task with official temporal splits; compare a GNN with a tabular baseline and inspect which joins would leak future information.\par
\textbf{Level.} Advanced\quad\href{https://papers.nips.cc/paper_files/paper/2024/hash/25cd345233c65fac1fec0ce61d0f7836-Abstract-Datasets_and_Benchmarks_Track.html}{Primary source $\nearrow$}\par
\end{minipage}\par\vspace{10pt}
\begin{minipage}{\linewidth}
{\large\bfseries P46 Robustness in Text-Attributed Graph Learning: Insights, Trade-offs, and New Defenses}\par
Runlin Lei, Lu Yi, Mingguo He, Pengyu Qiu, Zhewei Wei, Yongchao Liu, Chuntao Hong. ICLR 2026, 2026.\par
\textcolor{ink}{\small RECENT RESEARCH / 2026}\par
\textbf{Why read it.} Examines textual and structural perturbations in a common evaluation framework, connecting graph-learning robustness with the growing use of language-model features.\par
\textbf{Build or investigate.} On a public benchmark, compare text-only, edge-only and combined perturbations under fixed budgets; evaluate a GNN and a text encoder while separating poisoning from evasion.\par
\textbf{Level.} Advanced\quad\href{https://proceedings.iclr.cc/paper_files/paper/2026/hash/5133c31f29df92e056434cd86805cb6a-Abstract-Conference.html}{Primary source $\nearrow$}\par
\end{minipage}\par\vspace{10pt}
\Needspace{20\baselineskip}\section*{Twenty illustrative research projects}
A rigorous replication with a meaningful extension, a careful negative result, or a focused new method can all make a strong project. Publication acceptance is not a grading criterion. The required paper is submitted to the course; external conference or journal submission is optional and requires agreement among the authors.\par
These are scoped research suggestions. Agree on a dataset, question, and compute budget before beginning. Each project should report competitive and simple baselines, valid splits, ablations, error analysis, and reproducible artifacts.\par
\begin{minipage}{\linewidth}\subsection*{R01 When do biased walks help?}
Which graph structures make node2vec's walk bias useful?\par\textbf{Data.} Synthetic stochastic block models and Cora. \href{https://pytorch-geometric.readthedocs.io/en/latest/generated/torch_geometric.datasets.Planetoid.html}{Dataset documentation}\par
\textbf{Baselines.} DeepWalk, node2vec, degree features plus logistic regression\par
\textbf{Evaluation.} Node macro-F1 and accuracy; sweep homophily, walk length, p and q under fixed embedding budgets.\par
\textbf{Extension.} Relate performance to mixing time and community strength.\par
\textbf{Compute.} CPU-friendly\par
\end{minipage}\par\vspace{8pt}
\begin{minipage}{\linewidth}\subsection*{R02 PageRank as a learning prior}
Can personalized diffusion improve classification when labels are scarce?\par\textbf{Data.} Cora and CiteSeer. \href{https://pytorch-geometric.readthedocs.io/en/latest/generated/torch_geometric.datasets.Planetoid.html}{Dataset documentation}\par
\textbf{Baselines.} MLP, GCN, personalized PageRank propagation\par
\textbf{Evaluation.} Accuracy versus labeled fraction; fix splits and compare propagation depth and damping.\par
\textbf{Extension.} Choose damping using local graph structure without test labels.\par
\textbf{Compute.} CPU or small GPU\par
\end{minipage}\par\vspace{8pt}
\begin{minipage}{\linewidth}\subsection*{R03 Attention under noisy edges}
Does learned graph attention remain useful as irrelevant edges increase?\par\textbf{Data.} Cora plus controlled random edge additions. \href{https://pytorch-geometric.readthedocs.io/en/latest/generated/torch_geometric.datasets.Planetoid.html}{Dataset documentation}\par
\textbf{Baselines.} GCN, GraphSAGE, GAT\par
\textbf{Evaluation.} Accuracy and calibration versus edge corruption; match parameter counts and tuning budgets.\par
\textbf{Extension.} Compare attention entropy with robustness and test a sparsity regularizer.\par
\textbf{Compute.} Small GPU\par
\end{minipage}\par\vspace{8pt}
\begin{minipage}{\linewidth}\subsection*{R04 Learning across unlike neighbors}
When do neighborhood aggregation models fail on heterophilous graphs?\par\textbf{Data.} Roman-empire and Minesweeper benchmark splits. \href{https://pytorch-geometric.readthedocs.io/en/latest/generated/torch_geometric.datasets.HeterophilousGraphDataset.html}{Dataset documentation}\par
\textbf{Baselines.} MLP, GCN, GAT and a heterophily-aware model\par
\textbf{Evaluation.} Dataset-appropriate accuracy or ROC-AUC; preserve supplied splits and report multiple seeds.\par
\textbf{Extension.} Adapt the balance of self features and neighbor features.\par
\textbf{Compute.} Small GPU\par
\end{minipage}\par\vspace{8pt}
\begin{minipage}{\linewidth}\subsection*{R05 How deep can a GNN go?}
Which changes preserve predictive information as network depth increases?\par\textbf{Data.} Cora and synthetic community graphs. \href{https://pytorch-geometric.readthedocs.io/en/latest/generated/torch_geometric.datasets.Planetoid.html}{Dataset documentation}\par
\textbf{Baselines.} Plain GCN, residual GCN and normalization variants\par
\textbf{Evaluation.} Accuracy, embedding similarity and Dirichlet energy across depth; control training budget.\par
\textbf{Extension.} Separate optimization failure from oversmoothing with controlled interventions.\par
\textbf{Compute.} CPU or small GPU\par
\end{minipage}\par\vspace{8pt}
\begin{minipage}{\linewidth}\subsection*{R06 Rewiring graph bottlenecks}
Can a few added edges improve long-range information transfer?\par\textbf{Data.} Synthetic trees, barbell graphs and distance-label tasks. \href{https://networkx.org/documentation/stable/reference/generators.html}{Dataset documentation}\par
\textbf{Baselines.} GCN or MPNN with no rewiring, random rewiring and a topology-aware rule\par
\textbf{Evaluation.} Long-range task accuracy, path length and runtime at equal edge budgets.\par
\textbf{Extension.} Study the trade-off between shortcut benefits and structural distortion.\par
\textbf{Compute.} CPU-friendly\par
\end{minipage}\par\vspace{8pt}
\begin{minipage}{\linewidth}\subsection*{R07 A fair transformer comparison}
When does global attention earn its extra compute?\par\textbf{Data.} ZINC benchmark graphs. \href{https://pytorch-geometric.readthedocs.io/en/latest/generated/torch_geometric.datasets.ZINC.html}{Dataset documentation}\par
\textbf{Baselines.} GIN, a graph transformer and a local/global hybrid\par
\textbf{Evaluation.} Regression MAE, peak memory and runtime; match parameters and training budget.\par
\textbf{Extension.} Identify graph sizes or structures where global attention helps.\par
\textbf{Compute.} Small GPU; use a declared subset for pilots\par
\end{minipage}\par\vspace{8pt}
\begin{minipage}{\linewidth}\subsection*{R08 Which positional encoding matters?}
How stable are graph transformers to positional encoding choices?\par\textbf{Data.} ZINC and synthetic symmetric graphs. \href{https://pytorch-geometric.readthedocs.io/en/latest/generated/torch_geometric.datasets.ZINC.html}{Dataset documentation}\par
\textbf{Baselines.} Transformer without positional encoding, Laplacian encoding and random-walk encoding\par
\textbf{Evaluation.} MAE and permutation consistency; vary signs, eigenspaces and graph size.\par
\textbf{Extension.} Test sign-invariant or noise-robust encoding variants.\par
\textbf{Compute.} Small GPU\par
\end{minipage}\par\vspace{8pt}
\begin{minipage}{\linewidth}\subsection*{R09 Self-supervision with few labels}
Which graph augmentations help when supervision is limited?\par\textbf{Data.} Cora and CiteSeer. \href{https://pytorch-geometric.readthedocs.io/en/latest/generated/torch_geometric.datasets.Planetoid.html}{Dataset documentation}\par
\textbf{Baselines.} Supervised GCN, DGI and a contrastive graph encoder\par
\textbf{Evaluation.} Linear-probe accuracy across label budgets; tune augmentations on validation data only.\par
\textbf{Extension.} Replace uniform edge dropping with a structure-aware rule.\par
\textbf{Compute.} Small GPU\par
\end{minipage}\par\vspace{8pt}
\begin{minipage}{\linewidth}\subsection*{R10 Masked graph reconstruction}
Can reconstructing masked features improve robustness to missing attributes?\par\textbf{Data.} Cora and PubMed. \href{https://pytorch-geometric.readthedocs.io/en/latest/generated/torch_geometric.datasets.Planetoid.html}{Dataset documentation}\par
\textbf{Baselines.} MLP, GCN and GraphMAE-style pretraining\par
\textbf{Evaluation.} Classification accuracy and reconstruction loss across random and structured missingness.\par
\textbf{Extension.} Compare degree-aware masking with uniform masking.\par
\textbf{Compute.} Small GPU\par
\end{minipage}\par\vspace{8pt}
\begin{minipage}{\linewidth}\subsection*{R11 Predicting future collaborations}
How much do graph embeddings improve over simple link heuristics?\par\textbf{Data.} ogbl-collab. \href{https://ogb.stanford.edu/docs/linkprop/#ogbl-collab}{Dataset documentation}\par
\textbf{Baselines.} Common neighbors, Adamic-Adar, node2vec and GraphSAGE\par
\textbf{Evaluation.} Official Hits@50 with official temporal splits and negative candidates; document allowed edges.\par
\textbf{Extension.} Add recency-aware edge weights and audit temporal leakage.\par
\textbf{Compute.} Moderate GPU; prototype on a declared subset\par
\end{minipage}\par\vspace{8pt}
\begin{minipage}{\linewidth}\subsection*{R12 Learning from event streams}
When does temporal memory beat a static graph encoder?\par\textbf{Data.} Wikipedia temporal interaction data distributed with TGN. \href{https://github.com/twitter-research/tgn}{Dataset documentation}\par
\textbf{Baselines.} Static GraphSAGE, a recency heuristic and TGN\par
\textbf{Evaluation.} Average precision and ROC-AUC with time-ordered splits and explicit negative sampling.\par
\textbf{Extension.} Test memory resets, time encodings and unseen-node performance.\par
\textbf{Compute.} Small GPU\par
\end{minipage}\par\vspace{8pt}
\begin{minipage}{\linewidth}\subsection*{R13 Relations in a knowledge graph}
Do relation-aware messages improve completion beyond embedding models?\par\textbf{Data.} FB15k-237. \href{https://pytorch-geometric.readthedocs.io/en/latest/generated/torch_geometric.datasets.RelLinkPredDataset.html}{Dataset documentation}\par
\textbf{Baselines.} DistMult, TransE and R-GCN with a common scoring head where appropriate\par
\textbf{Evaluation.} Filtered MRR and Hits@10; prevent inverse-relation leakage and report relation-wise errors.\par
\textbf{Extension.} Allocate relation parameters by frequency or sparsity.\par
\textbf{Compute.} Small to moderate GPU\par
\end{minipage}\par\vspace{8pt}
\begin{minipage}{\linewidth}\subsection*{R14 Graph recommendations beyond popularity}
Can graph propagation improve recommendation for low-activity users?\par\textbf{Data.} MovieLens 100K stable benchmark. \href{https://grouplens.org/datasets/movielens/100k/}{Dataset documentation}\par
\textbf{Baselines.} Popularity, matrix factorization and LightGCN\par
\textbf{Evaluation.} Recall@20 and NDCG@20 on a time-ordered split; keep candidate sampling fixed.\par
\textbf{Extension.} Analyze accuracy and coverage across user-activity groups without inferring sensitive attributes.\par
\textbf{Compute.} CPU or small GPU\par
\end{minipage}\par\vspace{8pt}
\begin{minipage}{\linewidth}\subsection*{R15 Molecular learning under distribution shift}
Do stronger graph encoders generalize to new molecular scaffolds?\par\textbf{Data.} ogbg-molhiv. \href{https://ogb.stanford.edu/docs/graphprop/#ogbg-molhiv}{Dataset documentation}\par
\textbf{Baselines.} Fingerprint classifier, GCN and GIN\par
\textbf{Evaluation.} Official ROC-AUC on scaffold splits; add calibration and performance by scaffold frequency.\par
\textbf{Extension.} Compare pretraining or edge-feature treatments under equal tuning budgets.\par
\textbf{Compute.} Small GPU\par
\end{minipage}\par\vspace{8pt}
\begin{minipage}{\linewidth}\subsection*{R16 Generating diverse small graphs}
Which generative model best preserves structure without memorization?\par\textbf{Data.} Synthetic community and grid graphs. \href{https://networkx.org/documentation/stable/reference/generators.html}{Dataset documentation}\par
\textbf{Baselines.} Erdos-Renyi or block-model generator, GraphRNN and a small graph GAN\par
\textbf{Evaluation.} Validity, uniqueness, novelty, and degree/clustering distribution distances on held-out graphs.\par
\textbf{Extension.} Diagnose mode collapse as graph size and training diversity change.\par
\textbf{Compute.} Small GPU\par
\end{minipage}\par\vspace{8pt}
\begin{minipage}{\linewidth}\subsection*{R17 Adversarial versus diffusion generation}
How do adversarial and diffusion objectives trade quality for cost?\par\textbf{Data.} Small synthetic graphs; optional QM9 extension. \href{https://pytorch-geometric.readthedocs.io/en/latest/generated/torch_geometric.datasets.QM9.html}{Dataset documentation}\par
\textbf{Baselines.} A MolGAN-inspired discrete generator and a small DiGress-style model\par
\textbf{Evaluation.} Validity, diversity, novelty, structural fidelity and training/sample time at matched graph sizes.\par
\textbf{Extension.} If using molecules, add chemical-validity checks and compare conditional generation.\par
\textbf{Compute.} Moderate GPU; synthetic scope recommended\par
\end{minipage}\par\vspace{8pt}
\begin{minipage}{\linewidth}\subsection*{R18 Are graph explanations stable?}
Does a model explanation persist across seeds and harmless perturbations?\par\textbf{Data.} Synthetic motif graphs with known generating motifs. \href{https://pytorch-geometric.readthedocs.io/en/latest/generated/torch_geometric.datasets.BAShapes.html}{Dataset documentation}\par
\textbf{Baselines.} Gradient saliency, random edge masks and GNNExplainer\par
\textbf{Evaluation.} Fidelity, sparsity, motif precision/recall and agreement across repeated runs.\par
\textbf{Extension.} Build a stability-aware explanation selection rule; distinguish explanation from causality.\par
\textbf{Compute.} CPU or small GPU\par
\end{minipage}\par\vspace{8pt}
\begin{minipage}{\linewidth}\subsection*{R19 Robustness on a fixed edge budget}
Which defenses withstand targeted structural perturbations?\par\textbf{Data.} Cora or CiteSeer public benchmark graphs. \href{https://pytorch-geometric.readthedocs.io/en/latest/generated/torch_geometric.datasets.Planetoid.html}{Dataset documentation}\par
\textbf{Baselines.} GCN, GAT, edge-drop training and a simple graph denoiser\par
\textbf{Evaluation.} Clean and perturbed accuracy across equal budgets; distinguish random from targeted perturbations.\par
\textbf{Extension.} Test whether defense gains persist on unseen perturbation strategies in an offline benchmark.\par
\textbf{Compute.} Small GPU\par
\end{minipage}\par\vspace{8pt}
\begin{minipage}{\linewidth}\subsection*{R20 Do benchmark rankings survive protocol changes?}
How sensitive are model rankings to splits, tuning and random seeds?\par\textbf{Data.} Cora plus ogbn-arxiv, with a smaller scoped pilot if needed. \href{https://ogb.stanford.edu/docs/nodeprop/#ogbn-arxiv}{Dataset documentation}\par
\textbf{Baselines.} MLP, GCN, GraphSAGE and GAT\par
\textbf{Evaluation.} Accuracy distributions, rank stability, runtime and tuning cost; preserve official test protocols.\par
\textbf{Extension.} Measure how small search budgets alter conclusions and publish a reproducibility checklist.\par
\textbf{Compute.} Small to moderate GPU\par
\end{minipage}\par\vspace{8pt}
\end{document}